\documentclass[10pt]{article}
% Specify the margins and text size.
\setlength{\textwidth}{6.4in}
\setlength{\textheight}{9.5in}
\setlength{\oddsidemargin}{0pt}
\setlength{\evensidemargin}{0pt}
\setlength{\topmargin}{0pt}
\setlength{\hoffset}{.05in}
\setlength{\voffset}{-1in}

\setlength{\parskip}{5pt}
\setlength{\parindent}{0pt}

% Load some fonts and symbol packages
\usepackage{latexsym}
\usepackage{pifont}       % contains 'star' symbol for counterinsurgency handbook title
\usepackage{yfonts} 
\usepackage{amsmath}
\usepackage{amsfonts}

\usepackage{graphicx}     % actually, this is loaded by pstricks
\usepackage[T1]{fontenc}
\usepackage{ifthen}
\usepackage{pstricks,pst-grad,pst-text,pst-node,multido,pst-plot,calc,pst-3dplot}
%\usepackage[all]{xy}
%\usepackage{animate}

% The hyperref package inserts links into the pdf.
\definecolor{MyLinkColor}{rgb}{.1,.2,1}
\definecolor{MyCiteColor}{rgb}{.1,1,.2}
\definecolor{MyURLColor}{rgb}{.4,.4,.4}
\usepackage[backref=true,pagebackref=false,hyperindex,colorlinks=true,
  linkcolor=MyLinkColor,urlcolor=MyURLColor]{hyperref}


% The tweaklist package is something I found on the web.  It provides a simple interface
% for making changes to spacing used in the itemize and enumerate environments.  Comment
% this out if you don't care to use tweaklist.
\usepackage{tweaklist}
\renewcommand{\itemhook}{\setlength{\parskip}{2pt}\setlength{\parsep}%
{1pt}\setlength{\topsep}{0pt}\setlength{\itemsep}{0pt}}

\newcommand{\U}{\underline{\hspace{5pt}}}

\usepackage{listings}
\newcommand{\Z}{\hphantom{0}}

\begin{document}
\pagestyle{empty}
\lstset{language=R, showspaces=false, showstringspaces=false}

\href{http://www.su.edu}{\includegraphics[height=1.75cm]{sulogo.eps}}
\vspace{-1.79cm}

{{\ }\hfill\small
\begin{tabular}{cl}
& Math 207\\
& Introduction to Statistics\\
\end{tabular}
}
\setlength{\baselineskip}{1.05\baselineskip}
\medskip

\begin{center}
\textbf{\large  Worksheet:  The $\chi^2$ Test}
\end{center}
\smallskip



1.  Plotting the $\chi^2$ curve.

\hspace{20pt} a) Use the following to plot a $\chi^2$ curve with
  5 degrees of freedom.  Make a sketch the result.

\hspace{40pt}  \lstinline!>x <- seq(0, 30, by=0.2)!

\hspace{40pt}  \lstinline!>plot(x, dchisq(x, df=5), type='l', ylim=c(0, 0.15))!
\vspace{1.5in}

\hspace{20pt} b) Repeat part a) for a $\chi^2$ curve with
  10 degrees of freedom.  Make a sketch of the result using the\vspace{-4pt}

\hspace{20pt}\hphantom{b) } same axes you
  used in a).

\hspace{20pt} c) Repeat part a) for a $\chi^2$ curve with
  15 degrees of freedom. Make a sketch of the result using the\vspace{-4pt}

\hspace{20pt}\hphantom{b) } same axes you
  used in a).

\hspace{20pt} d) Describe the effect that increasing the number of degrees of
  freedom has on the shape of the curve.
\vspace{.5in}

\hspace{20pt} e) (Trick question) Estimate the area under each curve.
\bigskip

2. Find the area under the $\chi^2$-curve with 5 degrees of freedom to the right of
the given value. (Hint:  \lstinline!1-pchisq(x)!)

\hspace{20pt} a) 1.61\hspace{40pt} b) 9.24 \hspace{40pt} c) 15.09
\vspace{.25in}

3. Repeat exercise \#1 for a $\chi^2$-curve with 10 degrees of freedom.
\vspace{.4in}

4. A die is rolled 60 times.  The following frequency table is calculated.
Is the die fair? 

\begin{tabular}{|c|c|}\hline
      & observed\\
value & frequency\\\hline
1 & \Z5\\
2 & \Z7\\
3 & 17\\
4 & 16\\
5 & \Z8\\
6 & \Z7\\\hline
\end{tabular}

\vfill
\eject
{\ }

5. A die is rolled 60 times.  The following frequency table is calculated.
Is the die fair? 

\begin{tabular}{|c|c|}\hline
      & observed\\
value & frequency\\\hline
1 & \Z9\\
2 & 11\\
3 & 10\\
4 & \Z8\\
5 & 12\\
6 & 10\\\hline
\end{tabular}
\bigskip
\bigskip

6. A die is rolled 600 times.  The following frequency table is calculated.
Is the die fair? 

\begin{tabular}{|c|c|}\hline
      & observed\\
value & frequency\\\hline
1 & \Z90\\
2 & 110\\
3 & 100\\
4 & \Z80\\
5 & 120\\
6 & 100\\\hline
\end{tabular}
\bigskip
\bigskip

7. A die is rolled 60,000 times.  The following frequency table is calculated.
Is the die fair? 

\begin{tabular}{|c|c|}\hline
      & observed\\
value & frequency\\\hline
1 & 10,287\\
2 & 10,056\\
3 & \Z9,708\\
4 & 10,080\\
5 & \Z9,935\\
6 & \Z9,934\\\hline
\end{tabular}
\bigskip

8. A four-sided die is rolled several times.  
The following frequency table is calculated. Is the die fair? 

\begin{tabular}{|c|c|}\hline
      & observed\\
value & frequency\\\hline
1 & 105\\
2 & \Z90\\
3 & 111\\
4 & \Z94\\\hline
\end{tabular}

\vfill
Solutions: 1. d) Increasing the degrees of freedom shifts area to the right and 
   flattens the curve. e) $\mbox{area}=1$ \hspace{20pt} 
2. a) $90\%$ b) $10\%$ c) $1\%$ \hspace{20pt} 
3. a) $100\%$ b) $51\%$ c) $13\%$ \hspace{20pt} 
4. Probably not: $\chi^2=13.2$ and $P=2.2\%$\\ % \hspace{20pt}
5. Yes: $\chi^2=1.0$ and $P=96.3\%$ \hspace{20pt}
6. Maybe: $\chi^2=10.0$ and $P=7.5\%$ \hspace{20pt}\\
7. No: $\chi^2=18.6$ and $P=0.2\%$ \hspace{20pt}
8. Yes: $\chi^2=2.82$ and $P=42.0\%$ \hspace{20pt}
%9.  No: $\chi^2=152$ and $P\approx 0\%$ \hspace{20pt} 
%10. No: $\chi^2=6.58$ and $P=1.0\%$
\eject
%%%%%%%%%%%%%%%%%%%%%%%%%%%%%%%%%%%%%%

\href{http://www.su.edu}{\includegraphics[height=1.75cm]{sulogo.eps}}
\vspace{-1.79cm}

{{\ }\hfill\small
\begin{tabular}{cl}
& Math 207\\
& Introduction to Statistics\\
\end{tabular}
}

\medskip

\begin{center}
\textbf{\large  Worksheet:  The $\chi^2$ Test}
\end{center}
\smallskip



1.  Plotting the $\chi^2$ curve.

\hspace{20pt} a) Use the following to plot a $\chi^2$ curve with
  5 degrees of freedom.  Make a sketch the result.

\hspace{40pt}  \lstinline!>x <- seq(0, 30, by=0.2)!

\hspace{40pt}  \lstinline!>plot(x, dchisq(x, df=5), type='l', ylim=c(0, 0.15))!

{\color{blue}  There are some nice examples here:
\href{https://en.wikipedia.org/wiki/Chi-squared_distribution}{https://en.wikipedia.org/wiki/Chi-squared\U distribution}}
\vspace{1.2in}

\hspace{20pt} b) Repeat part a) for a $\chi^2$ curve with
  10 degrees of freedom.  Make a sketch of the result using the\vspace{-4pt}

\hspace{20pt}\hphantom{b) } same axes you
  used in a).

\hspace{20pt} c) Repeat part a) for a $\chi^2$ curve with
  15 degrees of freedom. Make a sketch of the result using the\vspace{-4pt}

\hspace{20pt}\hphantom{b) } same axes you
  used in a).

\hspace{20pt} d) Describe the effect that increasing the number of degrees of
  freedom has on the shape of the curve.
{\color{blue} It shifts the peak to the right.}
\vspace{.5in}

\hspace{20pt} e) (Trick question) Estimate the area under each curve. {\color{blue}The area under
any histogram is 1.}
\bigskip

2. Find the area under the $\chi^2$-curve with 5 degrees of freedom to the right of
the given value. (Hint:  \lstinline!1-pchisq(x)!)

\hspace{20pt} a) 1.61\hspace{40pt} b) 9.24 \hspace{40pt} c) 15.09

\hspace{20pt} a) {\color{blue}90.0\%}\hspace{40pt} b) {\color{blue}9.99\%} \hspace{40pt} c) {\color{blue}0.999\%}
\vspace{.25in}

3. Repeat exercise \#1 for a $\chi^2$-curve with 10 degrees of freedom.

\hspace{20pt} a) {\color{blue}99.9\%}\hspace{40pt} b) {\color{blue}50.9\%} \hspace{40pt} c) {\color{blue}12.9\%}
\vspace{.4in}

4. A die is rolled 60 times.  The following frequency table is calculated.
Is the die fair? 

{\color{blue}
  For a fair die, the expected frequence is $60\times \frac{1}{6} = 10$.  So,\\ 
$\chi^2 = \frac{(5-10)^2}{10} + \frac{(7-10)^2}{10} + \frac{(17-10)^2}{10} +
     \frac{(16-10)^2}{10} + \frac{(8-10)^2}{10} + \frac{(7-10)^2}{10}$.\\
The sum is $13.2$.
  So, the $P$-value is $1-\mbox{pchisq}(13.2, \mbox{df}=5)$\\ 
which is $2.16\mbox{\%}$.  
This supports the alternative hypothesis:  the die is not fair.}

\rput(14,0.8){\begin{tabular}{|c|c|}\hline
      & observed\\
value & frequency\\\hline
1 & \Z5\\
2 & \Z7\\
3 & 17\\
4 & 16\\
5 & \Z8\\
6 & \Z7\\\hline
\end{tabular}}

\vfill
\eject
{\ }

5. A die is rolled 60 times.  The following frequency table is calculated.
Is the die fair? 

{\color{blue}
  For a fair die, the expected frequence is $60\times \frac{1}{6} = 10$.  So,\\ 
$\chi^2 = \frac{(9-10)^2}{10} + \frac{(11-10)^2}{10} + \frac{(10-10)^2}{10} +
     \frac{(8-10)^2}{10} + \frac{(12-10)^2}{10} + \frac{(10-10)^2}{10}$.\\
The sum is $1.0$.
  So, the $P$-value is $1-\mbox{pchisq}(1.0, \mbox{df}=5)$\\ 
which is $96.3\mbox{\%}$.  
This supports the null hypothesis:  the die is fair.}

\rput(14, 0.8){\begin{tabular}{|c|c|}\hline
      & observed\\
value & frequency\\\hline
1 & \Z9\\
2 & 11\\
3 & 10\\
4 & \Z8\\
5 & 12\\
6 & 10\\\hline
\end{tabular}}
\bigskip
\bigskip

6. A die is rolled 600 times.  The following frequency table is calculated.
Is the die fair? 

{\color{blue}
  For a fair die, the expected frequence is $600\times \frac{1}{6} = 100$.  So,\\ 
$\chi^2 = \frac{(90-100)^2}{100} + \frac{(110-100)^2}{100} + \frac{(100-100)^2}{100} +
     \frac{(80-100)^2}{100} + \frac{(120-100)^2}{100} + \frac{(100-10)^2}{100}$.\\
The sum is $10.0$.
  So, the $P$-value is $1-\mbox{pchisq}(10.0, \mbox{df}=5)$\\ 
which is $7.5\mbox{\%}$.  
This weakly supports the  null hypothesis:  the die might be fair.}

\rput(14, 0.8){\begin{tabular}{|c|c|}\hline
      & observed\\
value & frequency\\\hline
1 & \Z90\\
2 & 110\\
3 & 100\\
4 & \Z80\\
5 & 120\\
6 & 100\\\hline
\end{tabular}}
\bigskip
\bigskip

7. A die is rolled 60,000 times.  The following frequency table is calculated.
Is the die fair? 

{\color{blue}
  For a fair die, the expected frequence is $60,\!000\times \frac{1}{6} = 10,\!000$.  So,\\ 
$\chi^2 = \frac{(10,287-10,000)^2}{10,000} + \frac{(10,056-10,000)^2}{10,000} + \cdots
      + \frac{(9,934-10,000)^2}{10,000}$.\\
The sum is $18.6$.
  So, the $P$-value is $1-\mbox{pchisq}(18.6, \mbox{df}=5)$\\ 
which is $0.2\mbox{\%}$.  
This  supports the  alternative hypothesis:  the die is not fair.}


\rput(14,0.8){\begin{tabular}{|c|c|}\hline
      & observed\\
value & frequency\\\hline
1 & 10,287\\
2 & 10,056\\
3 & \Z9,708\\
4 & 10,080\\
5 & \Z9,935\\
6 & \Z9,934\\\hline
\end{tabular}}
\bigskip\bigskip

8. A four-sided die is rolled several times.  
The following frequency table is calculated. Is the die fair? 

{\color{blue}
 The die was rolled a total of 400 times.  
  For a fair die, the expected frequence is $400\times \frac{1}{4} = 100$.  So,\\ 
$\chi^2 = \frac{(105-100)^2}{100} + \frac{(90-100)^2}{100} + \frac{(111-100)^2}{100} + 
      + \frac{(94-100)^2}{100}$.\\
The sum is $2.82$.
  So, the $P$-value is $1-\mbox{pchisq}(2.82, \mbox{df}=3)$\\ 
which is $0.2\mbox{\%}$.  
This  supports the  null hypothesis:  the die is  fair.}

{\ }

\rput(14,0.8){\begin{tabular}{|c|c|}\hline
      & observed\\
value & frequency\\\hline
1 & 105\\
2 & \Z90\\
3 & 111\\
4 & \Z94\\\hline
\end{tabular}}
\vfill
\eject

\end{document}

{\ }

9. As part of a study on the selection of grand juries in Alameda county, the 
educational level of grand jurors was compared with the county distribution
tabulated below.  Could a simple random sample of 62 people from 
the county show a distribution of educational level so different from the 
county-wide one?

\begin{tabular}{|lcc|}\hline
Educational    &        & Number of\\
level          & County & jurors\\\hline
Elementary     & \Z28.4\% & \Z1\\
Secondary      & \Z48.5\% & 10\\
Some college   & \Z11.9\% & 16\\
College degree & \Z11.2\% & 35\\\hline
Total &        100.0\%  & 62\\\hline
\end{tabular}
\vspace{1in}

10. (Hypothetical.) In a certain town, there are about one million eligible voters.  A
simple random sample of size 10,000 was chosen, to study the relationship 
between gender and participation in the last election.  The results are tabulated below.
Are gender and voting rate independent?

\begin{tabular}{|lcc|}\hline
            & Men & Women\\\hline
Voted       & 2,792 & 3,591\\
Didn't Vote & 1,486 & 2,131\\\hline
\end{tabular}

\vfill

Solutions: 1. d) Increasing the degrees of freedom shifts area to the right and 
   flattens the curve. e) $\mbox{area}=1$ \hspace{20pt} 
2. a) $90\%$ b) $10\%$ c) $1\%$ \hspace{20pt} 
3. a) $100\%$ b) $51\%$ c) $13\%$ \hspace{20pt} 
4. Probably not: $\chi^2=13.2$ and $P=2.2\%$\\ % \hspace{20pt}
5. Yes: $\chi^2=1.0$ and $P=96.3\%$ \hspace{20pt}
6. Maybe: $\chi^2=10.0$ and $P=7.5\%$ \hspace{20pt}\\
7. No: $\chi^2=18.6$ and $P=0.2\%$ \hspace{20pt}
8. Yes: $\chi^2=2.82$ and $P=42.0\%$ \hspace{20pt}
9.  No: $\chi^2=152$ and $P\approx 0\%$ \hspace{20pt} 
10. No: $\chi^2=6.58$ and $P=1.0\%$
\eject
\end{document}

